% Part: computability
% Chapter: computability-theory
% Section: russells-paradox

\documentclass[../../../include/open-logic-section]{subfiles}

\begin{document}

\olfileid{cmp}{thy}{rus}
\olsection{رسل دے پیراڈوکس نال تقابل}

ایس حصے دیاں دلیلاں دا رسل دے پیراڈوکس نال موازنہ کرنا،
تے اوہناں دیاں سانجھاں تے فرق ویکھنا، فائدہ مند اے:
\begin{enumerate}
% PNB-CMP-014: اصل پہلی شق وچ X غیر متعین اے؛سیٹ S دی تعریف تے رسل دے ثبوت مطابق S ہی اپنے آپ نال تقابل وچ آؤندا اے۔
\item رسل دا پیراڈوکس: من لو $S = \Setabs{x}{x \notin x}$۔
فیر $S
  \in S$ اُدوں تے صرف اُدوں جدوں $S \notin S$،
  تے ایہ تضاد اے۔

  \emph{نتیجہ:} ایہو جہا کوئی سیٹ~$S$ نہیں۔ ``سارے سیٹاں
  دا سیٹ'' موجود منّنا سیٹ نظریے دے ہور مسلّماں نال
  غیر موافق اے۔

\item رسل دے پیراڈوکس دا اک بدلیا روپ: من لو $F$ اوہ ``فنکشن''
  اے جیہڑا سارے فنکشناں دے سیٹ توں $\{ 0, 1 \}$ تک اے،
  تے ایس دی تعریف ایویں اے:
  \[
  F(f) =
  \begin{cases}
    1 & \text{جے $f$، $f$ دے دائرۂ تعریف وچ ہووے، تے $f(f) = 0$} \\
    0 & \text{ہور صورت وچ}
  \end{cases}
  \]
  ایداں دی دلیل وکھاندی اے کہ $F(F) = 0$ اُدوں تے صرف
  اُدوں جدوں $F(F) = 1$، تے ایہ تضاد اے۔

  \emph{نتیجہ:} $F$ فنکشن نہیں۔ ``سارے فنکشناں دا سیٹ''
  فنکشن دا دائرۂ تعریف بنن لئی بہت وڈا اے۔

\item قطری دلیل: من لو $f_0$، $f_1$، \dots جزوی قابلِ حساب
  فنکشناں دی فہرست اے، تے $G \colon \Nat \to
  \{ 0, 1 \}$ دی تعریف ایویں اے:
  \[
  G(x) =
  \begin{cases}
    1 & \text{جے $f_x(x)\downarrow = 0$} \\
    0 & \text{ہور صورت وچ}
  \end{cases}
  \]
  جے $G$ قابلِ حساب اے، تاں ایہ فنکشن $f_k$ اے،
  کسے $k$ لئی۔ پر فیر $G(k) = 1$ اُدوں تے صرف اُدوں
  جدوں $G(k) = 0$، تے ایہ تضاد اے۔

  \emph{نتیجہ:} $G$ قابلِ حساب نہیں۔ یاد رکھو کہ سیٹ نظریے
  دے مسلّماں مطابق $G$ فیر وی اک فنکشن اے؛ایتھے کوئی
  پیراڈوکس نہیں، صرف فرق دی وضاحت اے۔
\end{enumerate}

جزوی فنکشناں، قابلِ حساب فنکشناں، جزوی قابلِ حساب فنکشناں
وغیرہ دی ایہ گل الجھا سکدی اے۔ $\Nat$ توں $\Nat$ تک
سارے جزوی فنکشناں دا سیٹ چیزاں دا اک وڈا اکٹھ اے۔
اوہناں وچوں کجھ کُل نیں، کجھ قابلِ حساب نیں، کجھ کُل
تے قابلِ حساب دونوں نیں، تے کجھ نہ کُل نیں نہ قابلِ حساب۔
چیتے رکھو کہ جدوں اسی صرف ``فنکشن'' آکھدے آں، تاں
عام طور اُتے اسی کُل فنکشن دی گل کر رہے ہوندے آں۔ ایس طرح
چار قسمّاں نیں:
\begin{enumerate}
\item قابلِ حساب فنکشن
\item جزوی قابلِ حساب فنکشن جیہڑے کُل نہیں
\item فنکشن جیہڑے قابلِ حساب نہیں
\item جزوی فنکشن جیہڑے نہ کُل نیں نہ قابلِ حساب
\end{enumerate}
ایہ فرق صاف کرن لئی $\Nat$ توں $\Nat$ تک سارے جزوی
فنکشناں دا اک وڈا چوکور بنا کے ویکھو، تے اوہدے اندر
دو اک دوجے نال سانجھ رکھدے حصے نشان لاؤ: اک کُل
فنکشناں لئی، دوجا جزوی قابلِ حساب فنکشناں لئی۔
ایس خاکے توں بنن والے ہر حصے وچوں اک اک شے دی مثال
دے سکنا چنگی مشق اے۔

\end{document}
